\section{设计账本：构建多模态模型时真正要决定什么}

回到整讲，architecture 名称很多，但工程决策可以压缩成一张账本：

\begin{longtable}{L{0.18\textwidth}L{0.34\textwidth}L{0.38\textwidth}}
\toprule
决策 & 典型选项 & 主要权衡 \\
\midrule
Visual representation & continuous CLIP/SigLIP features；discrete VQ tokens & semantics vs reconstruction，理解 vs 生成 \\
Token count & fixed compression；dynamic resolution；tiling；frame sampling & 细节、context、attention compute \\
Position & 1D order；2D/3D MRoPE；explicit timestamps & 空间/时间结构与实现复杂度 \\
Fusion & input concat；cross-attention；cross-layer fusion & 接口简单、深层交互、显存 \\
Training & freeze/unfreeze stages；easy-to-hard curriculum & 稳定性、compute、catastrophic forgetting \\
Loss balance & token average；sample average；modality normalization & 长视频是否压过文本/图像 \\
Generation & external diffusion；native image tokens & 模块质量 vs 统一建模 \\
\bottomrule
\end{longtable}

\begin{knowledgebox}{Nota de clase}
多模态模型最常被低估的是 data pipeline：resolution buckets、OCR quality、caption contamination、video sampling、synthetic instruction、task balance 和 safety filtering。Architecture 图只展示系统的一小部分。
\end{knowledgebox}

